\section{La Giunzione Naturale in SQL: NATURAL JOIN e USING}

A differenza dell'algebra relazionale dove l'omissione della condizione trasforma implicitamente un Theta-Join in un Natural Join ($\bowtie$), il linguaggio SQL richiede l'uso di parole chiave esplicite per attivare questo comportamento ed evitare prodotti cartesiani accidentali.

Esistono due approcci sintattici per implementare la fusione delle colonne omonime tipica della giunzione naturale:

\paragraph{L'operatore puro: NATURAL JOIN}
Questa è la traduzione letterale dell'operatore matematico $\bowtie$.

\begin{lstlisting}[language=SQL, basicstyle=\ttfamily\small, keywordstyle=\color{blue}\bfseries]
SELECT *
FROM tabella_1
NATURAL JOIN tabella_2;
\end{lstlisting}

Quando il motore SQL incontra \texttt{NATURAL JOIN}, esegue autonomamente due operazioni:
\begin{enumerate}
    \item \textbf{Ricerca degli attributi:} Ispeziona gli schemi delle due relazioni alla ricerca di tutte le colonne aventi \textbf{esattamente lo stesso nome}.
    \item \textbf{Condizione e Fusione:} Genera un'uguaglianza implicita su tali colonne (es. \texttt{t1.id = t2.id}) e, nel \textit{Result Set} finale, \textbf{fonde} le colonne omonime in una sola, evitando ridondanze.
\end{enumerate}

\begin{warningbox}[La fragilità strutturale del NATURAL JOIN]
Nonostante la sua eleganza matematica, l'uso di \texttt{NATURAL JOIN} in produzione è universalmente considerato un \textbf{anti-pattern ingegneristico}. 
Il motivo risiede nella sua dipendenza occulta e non controllata dai nomi delle colonne. Se mesi dopo la stesura del codice un amministratore aggiungesse a entrambe le tabelle una colonna di servizio generica (es. \texttt{note} o \texttt{data\_inserimento}), il DBMS modificherebbe silenziosamente la condizione di giunzione pretendendo l'uguaglianza anche su queste nuove colonne (\texttt{AND t1.note = t2.note}). Questo distruggerebbe la logica dell'interrogazione originando risultati errati o insiemi vuoti, senza lanciare alcun errore di compilazione.
\end{warningbox}

\paragraph{L'alternativa ingegneristica: La clausola USING}

Per ovviare alla pericolosità del \texttt{NATURAL JOIN} preservandone però il vantaggio sintattico (la fusione delle colonne omonime), lo standard SQL ha introdotto la clausola \textbf{\texttt{USING}}.

Questa clausola va a sostituire \texttt{ON} e permette al progettista di dichiarare esplicitamente su quali colonne omonime deve avvenire la giunzione, rendendo il codice immune a future modifiche dello schema.

\textit{Sintassi con USING:}
\begin{lstlisting}[language=SQL, basicstyle=\ttfamily\small, keywordstyle=\color{blue}\bfseries]
-- Supponendo che entrambe le tabelle abbiano un attributo 'matricola'
SELECT *
FROM studente
JOIN esame USING (matricola);
\end{lstlisting}

A differenza della classica clausola \texttt{ON} (che restituirebbe due colonne distinte: \texttt{studente.matricola} ed \texttt{esame.matricola}), la clausola \texttt{USING} istruisce il database a \textbf{fondere} l'attributo in una singola colonna, rispecchiando in totale sicurezza il comportamento algebrico della giunzione naturale.